\chapter{Deep Learning Architectures}

\section{Transition to Deep Learning}
While tabular models are efficient, they rely on hand-crafted features. Deep learning models can learn hierarchal representations directly from raw acoustic data (spectrograms), potentially capturing subtle anomalies that manual features might miss.

\section{Architecture: ResNet-34 Convolutional Neural Network}
Our primary deep learning experiment utilized a custom \textbf{ResNet-34} architecture. This model moves beyond the shallower ResNet-18 by employing a $(3, 4, 6, 3)$ residual block configuration, providing a deeper acoustic manifold for capturing complex mechanical signatures.

\subsection{Regularization and Training Dynamics}
To ensure robust generalization on the factory floor, we implemented a state-of-the-art regularization suite:
\begin{itemize}
    \item \textbf{SpecAugment:} We applied on-the-fly spectral masking, including Frequency Masking ($\nu=30$) and Time Masking ($\tau=80$). This prevents the model from overfitting to specific transient spectral artifacts.
    \item \textbf{Mixup Augmentation:} In 30\% of training iterations, we employed Mixup ($\alpha=0.2$). This encourages the model to learn linear decision boundaries between acoustic classes, improving its performance on boundary-case anomalies.
    \item \textbf{Label Smoothing:} We used a smoothing factor of $0.1$ in our Cross-Entropy loss to prevent overconfident predictions and improve calibration.
    \item \textbf{Weight Decay:} A L2 penalty of $0.01$ and a Dropout layer ($p=0.5$) before the final classification head were used to constrain model complexity.
\end{itemize}

\section{Architectural Explainability (XAI)}
A unique feature of our ResNet-34 implementation is its \textbf{Diagnostic Mode}. The model architecture is designed to return internal \textbf{Feature Maps} from every residual block during inference. By aggregating and visualizing these activations, engineers can see exactly which frequency bands or temporal impulses triggered an anomaly detection. This "Explainable AI" approach transforms the model from a black box into a transparent tool for root-cause analysis on the factory floor.

\subsection{Experimental Iterations}
We conducted several experimental runs with varying clip durations and dataset configurations to find the optimal training setup for the ResNet model.

\begin{table}[H]
    \centering
    \caption{ResNet Performance across Experimental Runs}
    \begin{tabular}{@{}lcccc@{}}
        \toprule
        \textbf{Run ID} & \textbf{Clip Duration} & \textbf{Dataset Source} & \textbf{Val Accuracy} & \textbf{Val F1-Score} \\ \midrule
        Run 1 & 10.0s & Raw & 96.79\% & 0.9515 \\
        Run 2 & 5.0s & Raw & 98.89\% & 0.9801 \\
        Run 3 & 5.0s & Mel-NPY & 96.94\% & 0.9438 \\
        Run 4 & 5.0s & Mel-NPY & 96.78\% & 0.9494 \\
        Run 5 & \textbf{5.0s} & \textbf{Raw (Final)} & \textbf{99.57\%} & \textbf{0.9928} \\ \bottomrule
    \end{tabular}
\end{table}

The final run (Run 5) demonstrated that using 5-second clips from the raw dataset provided the most robust representation for the ResNet backbone, likely due to the higher temporal resolution of the training windows.

\section{EfficientAT: Advanced Machine Fault Classification}
To push the state-of-the-art in machine sound classification, we developed \textbf{EfficientAT} (EfficientNet + Audio Frontend + Temporal Attention). This model moves beyond fixed feature extraction by integrating a differentiable audio processing pipeline directly into the neural network, allowing for end-to-end classification from raw \texttt{.wav} audio.

\subsection{Architectural Components}
EfficientAT is composed of three sequential stages designed to capture spectral, spatial, and temporal features:

\subsubsection{Component 1: Differentiable Audio Frontend}
The \texttt{AudioFrontend} module converts raw waveforms into a rich 3-channel representation within the computational graph, eliminating the need for offline preprocessing:
\begin{itemize}
    \item \textbf{Mel-Spectrogram:} Computes power spectrograms across 128 Mel bands.
    \item \textbf{Log Compression:} Converts power values to decibels, approximating human auditory perception.
    \item \textbf{Delta features ($\Delta$ and $\Delta\Delta$):} Computes first and second-order temporal derivatives to encode the velocity and acceleration of spectral changes.
    \item \textbf{Channel Stacking:} Concatenates these three maps to produce a 3-channel pseudo-image, making the audio data compatible with image-pretrained backbones.
\end{itemize}

\subsubsection{Component 2: EfficientNet-B0 Backbone}
We employ an \textbf{EfficientNet-B0} backbone for its high efficiency. The model is configured to preserve temporal structure:
\begin{itemize}
    \item \textbf{Spatial preservation:} Global pooling is disabled to retain a 2D feature map.
    \item \textbf{Frequency Pooling:} Frequency information is collapsed via adaptive average pooling, yielding a 1D sequence of 1,280-dimensional vectors along the time axis.
\end{itemize}

\subsubsection{Component 3: Temporal Attention (Transformer)}
Machine faults often manifest as rhythmic anomalies or transient impulses. To capture these long-range dependencies, we integrate a \textbf{Transformer Encoder}:
\begin{itemize}
    \item \textbf{Multi-Head Attention:} A 2-layer encoder with 4 attention heads allows the model to "attend" to specific moments in the 11-second recording that characterize a fault.
    \item \textbf{Global Context:} Unlike local convolutions, self-attention captures relationships across the entire temporal window, which is crucial for distinguishing steady-state machine noise from intermittent faults.
\end{itemize}

\subsection{Experimental Results and Performance}
We evaluated EfficientAT in two configurations to characterize its performance in real-world (Raw WAV) and ideal (Preprocessed) conditions.

\begin{table}[H]
    \centering
    \caption{EfficientAT Experimental Benchmarking (Cluster-Aware Split)}
    \begin{tabular}{@{}lccc@{}}
        \toprule
        \textbf{Configuration} & \textbf{Test Accuracy} & \textbf{Macro F1-Score} & \textbf{Inference Latency} \\ \midrule
        \textbf{Raw WAV (EfficientAT + CGAN)} & \textbf{89.38\%} & \textbf{0.8407} & \textbf{21.4 ms / sample} \\
        Baseline (Standard Split) & 99.82\% & 0.9972 & 21.4 ms / sample \\ \bottomrule
    \end{tabular}
\end{table}

The 21.4 ms latency for the end-to-end pipeline includes the full on-the-fly feature extraction, supporting a throughput of $\sim$47 samples/second on a single GPU. This demonstrates that EfficientAT is production-ready for real-time monitoring without the complexity of a separate signal processing pipeline.

\begin{figure}[H]
    \centering
    \begin{tikzpicture}[
        node distance=0.55cm and 0.4cm,
        box/.style={rectangle, rounded corners=4pt, draw=primary, fill=primary!5,
                    minimum width=3.8cm, minimum height=0.75cm,
                    font=\small\bfseries, align=center},
        arrow/.style={-{Stealth[length=5pt]}, thick, color=accent},
        label/.style={font=\scriptsize\color{gray}, align=center}
    ]
    \node[box] (wav) {Raw .wav\\(B, 1, 176000)};
    \node[box, below=of wav] (mel) {MelSpectrogram\\(B, 1, 128, 344)};
    \node[box, below=of mel] (db) {Log Compression\\AmplitudeToDB};
    \node[box, below=of db] (delta) {Delta Features\\$\Delta$ + $\Delta\Delta$};
    \node[box, below=of delta] (stack) {Channel Stack\\(B, 3, 128, 344)};
    \node[box, below=of stack, fill=accent!15, draw=accent] (eff) {EfficientNet-B0\\Feature Maps (B, 1280, 4, 11)};
    \node[box, below=of eff] (pool) {Freq. Pool\\(B, 1280, 11)};
    \node[box, below=of pool] (proj) {Conv1d Projection\\(B, 128, 11)};
    \node[box, below=of proj, fill=orange!15, draw=orange] (attn) {Transformer Encoder\\2 Layers, 4 Heads};
    \node[box, below=of attn] (clf) {Classifier Head\\Dropout + Linear};
    \node[box, below=of clf, fill=green!15, draw=green] (out) {6-Class Output\\M1-M3 Normal/Abnormal};
    
    \foreach \a/\b in {wav/mel, mel/db, db/delta, delta/stack, stack/eff,
                       eff/pool, pool/proj, proj/attn, attn/clf, clf/out}
      \draw[arrow] (\a) -- (\b);
    
    \node[label, right=0.25cm of eff] {EfficientNet-B0\\$\sim$5.3M params};
    \node[label, right=0.25cm of attn] {Long-range\\temporal modelling};
    \end{tikzpicture}
    \caption{Full EfficientAT Pipeline Architecture}
    \label{fig:efficientat_architecture}
\end{figure}

\section{SSLAM: Simple Spectral-Linear Attention Module}
To bridge the gap between high-complexity models like EfficientNet and low-latency requirements, we developed the \textbf{SSLAM} (Simple Spectral-Linear Attention Module). This architecture targets the "mid-range" of inference, providing higher representational power than a basic ResNet-18 without the heavy Transformer overhead of EfficientNet.

\subsection{Architecture Design}
The SSLAM architecture consists of:
\begin{itemize}
    \item \textbf{CNN Feature Extraction:} A lightweight 2-layer convolutional block with Batch Normalization and PReLU activations.
    \item \textbf{Multi-Head Attention (MHA):} Instead of a full Transformer, we employ a single 4-head attention layer directly on the flattened feature maps. This allows the model to learn global dependencies across frequency bands.
    \item \textbf{Dual-Head Classification:} A primary classification head for final predictions and an auxiliary head for intermediate loss calculation, improving gradient flow during training.
\end{itemize}

\section{Addressing Class Imbalance with Conditional GANs (CGAN)}
Industrial acoustic datasets are inherently imbalanced; mechanical failures are rare compared to steady-state operation. In our dataset, abnormal signatures for machines M1, M2, and M3 were significantly underrepresented. To address this without mere oversampling (which risks overfitting to specific recording sessions), we implemented a \textbf{Conditional Generative Adversarial Network (CGAN)} to synthesize realistic Mel-spectrograms.

\subsection{CGAN Architecture and Training}
The generative pipeline follows a minimax game between two competing networks:
\begin{itemize}
    \item \textbf{Generator ($G$):} Utilizes a 1D upsampling architecture (along the time axis) with residual blocks and transposed convolutions. It maps a 128-dimensional noise vector $\mathbf{z}$ and a class embedding $y$ to a synthetic $[1, 128, 344]$ Mel-spectrogram. The use of \textbf{Hinge Loss} ensures stable gradient flow during synthesis.
    \item \textbf{Discriminator ($D$):} A PatchGAN-style classifier that processes the spectrogram-class pairs. It is tasked with distinguishing real spectral manifolds from synthetic ones while ensuring the output aligns with the requested class condition.
\end{itemize}

\subsection{Two-Phase Training Strategy}
To maximize the utility of synthetic data, the \textbf{EfficientAT} model was trained in two distinct phases:
\begin{enumerate}
    \item \textbf{Phase 1 (Real Pre-training):} The model is trained solely on the real dataset using standard Cross-Entropy loss to learn the foundational acoustic manifold of the factory environment.
    \item \textbf{Phase 2 (Synthetic Fine-tuning):} We generate 400 high-fidelity samples for each minority abnormal class (M1-A, M2-A, M3-A). The model is then fine-tuned on the combined \texttt{Real + Synthetic} set using \textbf{Focal Loss} ($\gamma=2.0$) and \textbf{Hard Class Weights} ($w \in [2.5, 4.5]$) to force the network to focus on hard-to-classify abnormal signatures.
\end{enumerate}

\subsection{Asymmetric Label Smoothing}
A novel regularization technique, \textbf{Asymmetric Label Smoothing}, was employed during the fine-tuning phase. While majority classes received standard smoothing ($0.1$) to prevent overconfidence, minority abnormal classes received reduced smoothing ($0.05$). This preserves the sharpness of the minority class gradients, preventing the synthetic samples from being "diluted" by the dominant normal signatures.

\subsection{Rigorous Evaluation: Cluster-Aware Split}
Unlike early experiments that utilized standard block splitting (reaching $>99\%$ accuracy), the final evaluation was conducted using a \textbf{DBSCAN Cluster-Aware Split}. By grouping recordings based on their deep spectral embeddings before splitting, we ensured that the model could not "memorize" recording sessions. Under this strict regime, the \textbf{EfficientAT + CGAN} pipeline achieved a validated \textbf{89.38\% accuracy} and a \textbf{0.8407 Macro F1-score}, representing a highly generalized diagnostic capability that is robust against factory-floor noise variations.

% End of Deep Learning Architectures
